\section{Two certificate reductions and two measured alternatives}
\label{sec:architecture}

The principal improvements here are a fixed polynomial majorant for the
CRT bound and one-sided witnesses for every Pell congruence. The arithmetic
proofs are independent of the proof-calculus calibration. Their composition
with that calibration does not assume consistency.

\subsection{One fixed majorant instead of a second row circuit}

Let $E,F$ be the nonnegative arithmetic circuits splitting the row
polynomial. Take
\[
 R=4+N+x+y+m+M_0+U.
\]
This bounds each external row input, the row index, the labels bounded by
$m$, the accumulator values bounded by $y$, and the other local coordinates
bounded by $U$. It also bounds every range-certificate cutoff, including
$m+1,U+1,y$, and $N$.

Assign exponent $1$ to inputs; exponent $0$ to constants $0,1$; and exponent
$1$ to constants $2,3,4$. At addition take the larger input exponent plus
one, and at multiplication take their sum. Every circuit node with assigned
exponent $d$ is at most $R^d$. Indeed, multiplication preserves the
induction bound, and addition uses $2R^d\le R^{d+1}$ because $R\ge4$.

Let $d=\max(d_E,d_F)$, and choose a fixed power of two $\widehat d\ge d$.
Then
\[
 0\le E\le K,\qquad0\le F\le K,\qquad K=R^{\widehat d}.
\]
These are \emph{separate} bounds. They do not assert $E+F\le K$.
The CRT soundness argument only requires separate bounds: a decoded
congruence modulo a prime $p>K$ between $E,F\in[0,K]$ is equality.
A stronger optional convention replaces $d$ by $\max(d_E,d_F)+1$,
which also bounds $E+F$; the benchmark appendix uses that harmlessly
larger convention. Both audited source circuits fit under
$\widehat d=2^{23}$. Exactly 23 squaring gates compute $K$, without
introducing a variable-exponent power graph. Put $T=K!$ as before.
The radius also ensures that all range factorials divide $T$, so ordinary
CRT gives completeness of the factorial-free binomial certificates.

The choice is noncircular: proof and labels, carrier, obligation stream,
beta scale and code, local bound $U$, then $R,K,T$, then CRT representatives.
The earlier independent radius benchmark decreased 2799 natural variables
to 2672. Subsequent source factoring and forward-quotient substitutions in
the main construction improve this further; 2672 is not the final record.

\subsection{All indexed Pell congruences can point forward}

For $A\ge2$, write $X_A(j),Y_A(j)$ for the canonical Pell sequence, with
\[
\begin{split}
 X_A(0)&=1,\quad Y_A(0)=0,\\
 X_A(j+1)&=AX_A(j)+(A^2-1)Y_A(j),\\
 Y_A(j+1)&=X_A(j)+AY_A(j).
\end{split}
\]
Simultaneous induction proves monotonicity of both coordinates in $A$.
It also proves $Y_A(j)\ge j$.

The existing indexed-Pell completeness construction, for target index
$k>0$, chooses
\[
 X=X_A(k),\quad Y=Y_A(k),\quad h=2kY,\quad
 U=X_A(h),\quad V=Y_A(h).
\]
Here $h\ge2$, so $U>A$. It chooses $B$ simultaneously satisfying
$B\equiv A\pmod U$ and $B\equiv1\pmod{4Y}$, then takes
$S=X_B(k)$ and $R_1=Y_B(k)$. The latter index is the original target
$k$, not the auxiliary $h$. Since $0\le A<U$, the residue condition
forces $B\ge A$. Hence $S\ge X$ and $R_1\ge k$. The same completeness
witness therefore admits natural quotients satisfying
\[
 B=A+Uq_B,\qquad S=X+Uq_S,\qquad R_1=k+4Yq_R.
\]
Each equation implies its former congruence, so soundness is preserved
by restriction. The canonical simultaneous witnesses prove completeness;
there is no claim that every old auxiliary witness has the new orientation.

The exact source is Mathlib's \texttt{PellMatiyasevic.lean} at commit
\texttt{4dc43842a11f257ffba2a611c943a133a65cb4b2}, lines 761--807.
The inert source, its Apache license, exact blob identity and SHA256 are
included in the architecture packet. The construction uses $V>0$,
$Y^2\mid V$ and $\gcd(U,4Y)=1$ simultaneously, as required by the
existing indexed-Pell theorem.

\subsection{The power congruence can also point forward}

For $A\ge n\ge1$, define
\[
 F_j=X_A(j)-(A-n)Y_A(j).
\]
Then $F_0=1,F_1=n$ and $F_{j+1}=2AF_j-F_{j-1}$. We prove
$F_j\ge nF_{j-1}\ge0$ by induction. Its induction step is
\[
 F_{j+1}-nF_j=(2A-n)F_j-F_{j-1}
 \ge\bigl(n(2A-n)-1\bigr)F_{j-1}\ge0.
\]
Thus $F_k\ge n^k$. If $T=2An-n^2-1>0$, the recurrence also gives
$F_k\equiv n^k\pmod T$, since the recurrence error of $n^j$ is a
multiple of $T$. Therefore
\[
 X_A(k)+nY_A(k)=AY_A(k)+n^k+Tq
 \qquad(q\in\mathbb N).
\]
The positive power graph's last norm implies $A>W\ge n$:
$A^2=1+W^3(W+2)Z^2$, and $A\ge2$ forces $Z\ge1$.
Its canonical witnesses therefore satisfy the forward equation. Conversely,
forgetting the sign of the quotient recovers the inherited exact power
criterion. Zero base or exponent is not silently included: all calls in
this specialized compiler are separately proved positive.

\subsection{Negative CRT moduli as a valid alternative}

Suppose the row condition is $\prod_j C_j=0$ and $R_0=0$, where all
$C_j,R_0$ are sums of squares. Bound each $C_j$ and $R_0$ separately
by $K$, with $K\ge4,N$ and all range cutoffs. Put
\[
 T=K!,\quad D_i=(i+1)T-1,\quad B=\prod_{i<N}D_i.
\]
The $D_i$ are pairwise coprime by the usual difference argument, and each
is $3\pmod4$. Hence each has a prime divisor $p\equiv3\pmod4$;
since $p\nmid K!$, necessarily $p>K$.
Using the usual binomial range certificates, the index congruence
$B\mid(C+1)T-1$, and
\[
 \left(\prod_jC_j(C,A)\right)^2+R_0(C,A)^2\equiv0\pmod B
\]
recovers every row. Indeed, over such a prime, a sum of two squares is
zero only when both terms are zero. Primality then selects a vanishing
$C_j$, and the separate bounds turn the selected branch and $R_0$ into
exact zeros. Ordinary CRT proves completeness.

No variable-length product remains in the output. With
$Q=(1+(N+1)T)^N+2$ one has
\[
 B=\left(T^N N!\binom{Q+N}{N}\right)\bmod(TQ+1).
\]
The numerator is $\prod_{j=1}^N(TQ+jT)$, which reduces to
$\prod_{j=1}^N(jT-1)$, and $B<Q<TQ+1$. For $N=0$ the product is
$1$, the congruences modulo $1$ are vacuous, and the same identity holds.
The rounded-radius negative-modulus benchmark has 2678 natural variables,
six more than its matched positive-modulus comparison. It is a proved
alternative, not the recommended minimum.

\subsection{Factorial separation without a coefficient height}

If $M\ge1$ and $\max(m-1,0)!\mid M$, then the numbers
$1+jM$, $1\le j\le m$, are pairwise coprime. A common divisor for
labels $j<k$ is coprime to $M$ and divides $k-j$, while $k-j$ divides
$M$. Therefore it is $1$. Prime valuations now recover every positive
label multiplicity from a product equality, for arbitrary proof length.
The factors themselves need not be prime. Label zero contributes the
unit factor and remains invisible. For $m=0$ the target factor $1+M>1$
makes the equality impossible; for $m=1$ any positive $M$ works.

Thus $M=m!$ is a correct replacement carrier when $m\ge1$. It is
chosen before the obligation stream, never from its later CRT bound.
The matched native-graph benchmark costs 2872 instead of 2672 variables:
one extra factorial, one extra binomial and four extra power graphs.
A completely fixed $M$ without a label restriction is unsound, as shown by
$1+(M+2)M=(1+M)^2$ and the underivable single rule $(1,0)\to M+2$.

All arguments above are finite arithmetic theorems formalizable in ZFC.
They remain internal theorems in nonstandard models; they do not rely on
external standardness of a proposed certificate. The included finite tests
and literal polynomial replays verify implementation fidelity, not a new
kernel proof of those all-input statements.
